%
% jordan-ungleichung.tex -- Sinusbogen oberhalb seiner Sehne
%
% (c) 2026 Damien Flury, OST Ostschweizer Fachhochschule
%
\documentclass[tikz]{standalone}
\usepackage{amsmath}
\usepackage{times}
\usepackage{txfonts}
\usetikzlibrary{arrows,math}
\definecolor{darkred}{rgb}{0.8,0,0}
\definecolor{darkviolet}{rgb}{0.45,0,0.6}
\begin{document}
\def\skala{1}
\begin{tikzpicture}[>=latex,thick,scale=\skala,xscale=1.7,yscale=1.5]

\def\hpi{1.5707963}
\def\pipi{3.1415927}

% Flaeche zwischen Sehne und Sinusbogen, wo die Ungleichung gilt
\fill[black!8]
	plot[domain=0:\hpi,samples=100] ({\x},{sin(\x r)}) -- (0,0) -- cycle;

% Achsen
\draw[->] (-0.18,0) -- (\pipi+0.45,0) coordinate[label={$\theta$}];
\draw[->] (0,-0.18) -- (0,2.30) coordinate[label={right:$y$}];

% Trennlinie bei pi/2
\draw[color=gray,dashed,line width=0.6pt] (\hpi,0) -- (\hpi,1);

% Sinusbogen
\draw[color=darkred,line width=1.4pt]
	plot[domain=0:\pipi,samples=200] ({\x},{sin(\x r)});

% Sehne, ueber pi/2 hinaus weitergefuehrt
\draw[color=darkviolet,line width=1.4pt] (0,0) -- (\pipi,2);

% Markierungen auf den Achsen
\draw (\hpi,-0.05) -- (\hpi,0.05);
\node at (\hpi,0) [below] {$\frac{\pi}{2}$};
\draw (\pipi,-0.05) -- (\pipi,0.05);
\node at (\pipi,0) [below] {$\pi$};
\draw (-0.04,1) -- (0.04,1);
\node at (0,1) [above left] {$1$};
\draw (-0.04,2) -- (0.04,2);
\node at (0,2) [above left] {$2$};
\node at (0,0) [below left] {$0$};

% Beschriftung der Kurven
\node[color=darkred] at (2.45,0.638) [above right] {$y=\sin\theta$};
\node[color=darkviolet] at (0.72,0.458) [below right] {$y=\frac{2\theta}{\pi}$};

\end{tikzpicture}
\end{document}
